% Compile in terminal by running:
% latexmk -C CV.tex # Wipes all the cached build files
% latexmk -pdf CV.tex

\documentclass[10pt,a4paper]{article}

%----------------------------------------------------------------------------------
% Packages
%----------------------------------------------------------------------------------
\usepackage[a4paper,top=1.35cm,bottom=1.25cm,left=1.55cm,right=1.55cm]{geometry}
\usepackage[T1]{fontenc}
\usepackage[utf8]{inputenc}
\usepackage{lmodern}
\usepackage{microtype}
\usepackage{enumitem}
\usepackage{xcolor}
\usepackage[hidelinks]{hyperref}
\usepackage{titlesec}
\usepackage{tabularx}
\usepackage{ragged2e}

%----------------------------------------------------------------------------------
% General formatting
%----------------------------------------------------------------------------------
\pagestyle{empty}

\setlength{\parindent}{0pt}
\setlength{\parskip}{0pt}

\definecolor{headingcolor}{RGB}{25,55,90}
\definecolor{rulecolor}{RGB}{155,155,155}

% Section formatting
\titleformat{\section}
  {\large\bfseries\color{headingcolor}}
  {}
  {0pt}
  {}
  [\vspace{0.15em}\color{rulecolor}\titlerule]

\titlespacing*{\section}
  {0pt}{1em}{0.6em}

% Compact but breathable lists
\setlist[itemize]{
    leftmargin=1.25em,
    itemsep=0.2em,
    topsep=0.2em,
    parsep=0em,
    partopsep=0em,
    after=\vspace{0.4em} % Adds consistent space after lists
}

\setlist[enumerate]{
    leftmargin=1.45em,
    itemsep=0.4em,
    topsep=0.2em,
    parsep=0em,
    after=\vspace{0.4em}
}

%----------------------------------------------------------------------------------
% Custom commands
%----------------------------------------------------------------------------------

\newcommand{\entry}[4]{
    \noindent\textbf{#1} \hfill \textbf{#2}\\
    \textit{#3} \hfill #4\vspace{0.3em}\par
}

\newcommand{\project}[1]{
    \noindent\textbf{\textit{#1}}\vspace{0.1em}\par
}

\newcommand{\skill}[2]{
    \noindent\textbf{#1:} #2\vspace{0.25em}\par
}

%----------------------------------------------------------------------------------
% Document
%----------------------------------------------------------------------------------
\begin{document}

%==================================================================================
% HEADER
%==================================================================================

\begin{center}

{\fontsize{24}{28} \selectfont \bfseries Ankit Mukherjee, MSc}

\vspace{0.3em}

{\normalsize Computational Biologist}

\vspace{0.45em}

Barrackpore, Kolkata, West Bengal, India
\quad | \quad
+91 9330111538
\quad | \quad
\href{mailto:ankit0204.am@gmail.com}{ankit0204.am@gmail.com}

\vspace{0.2em}

\href{https://www.linkedin.com/in/ankit-mukherjee-647ab81ba/}{LinkedIn}
\quad | \quad
\href{https://github.com/AnkitMukherji}{GitHub}
\quad | \quad
\href{https://orcid.org/0009-0005-5847-5986}{ORCiD}

\end{center}

\vspace{0.4em}

%==================================================================================
% RESEARCH INTERESTS
%==================================================================================

\section{Research Interests}

Computational Genomics \textbullet\ Population Genomics \textbullet\ Human Genetic Variation
\textbullet\ Pharmacogenomics \textbullet\ Transcriptomics \textbullet\ Functional Genomics
\textbullet\ Disease Biology \textbullet\ Precision Medicine \textbullet\ Statistical Genomics

%==================================================================================
% EDUCATION
%==================================================================================

\section{Education}

\entry
{MSc in Bioinformatics (DBT-funded)}
{2022--2024}
{Savitribai Phule Pune University, Pune}
{CGPA: \textbf{9.93/10}}

\entry
{BSc in Microbiology}
{2019--2022}
{University of Calcutta, Kolkata}
{CGPA: \textbf{9.03/10}}

%==================================================================================
% RESEARCH EXPERIENCE
%==================================================================================

\section{Research Experience}

\entry
{Project Associate}
{2024--Present}
{CSIR--Institute of Genomics and Integrative Biology (CSIR--IGIB), Delhi}
{}

\project{Population Pharmacogenomics and Precision Medicine -- GenomeIndia Project}
\begin{itemize}
\item Investigating the landscape of \textbf{clinically actionable pharmacogenomic variation} across the Indian population using whole-genome sequencing data from \textbf{10,000 individuals representing 83 genetically diverse Indian populations}.
\item Curated and harmonized pharmacogenomic evidence from \textbf{PharmGKB, CPIC, PharmVar, and ClinPGx} to identify clinically actionable SNPs, indels, and star alleles relevant to drug response.
\item Characterized \textbf{population-specific allele-frequency variation} across Indian populations and compared pharmacogenomic profiles with global populations to identify differences with potential implications for precision medicine.
\item Analyzed clinically important pharmacogenes and star-allele haplotypes to infer \textbf{drug metabolizer phenotypes}, identifying population groups with potentially altered responses to commonly prescribed therapeutics.
\end{itemize}

\project{GenomeIndia Database (GI-DB): \href{https://ibdc.dbt.gov.in/gidb/}{ibdc.dbt.gov.in/gidb/}}
\begin{itemize}
\item Contributing author to \textbf{GI-DB}, a publicly accessible genomic database cataloguing genetic variation identified from \textbf{10,000 whole genomes representing 83 Indian populations}.
\item Contributed to large-scale \textbf{variant processing, quality assessment, normalization, functional annotation and population-level analyses} supporting the construction of a comprehensive Indian genomic variation resource.
\end{itemize}

\project{Blood Transcriptomics and Disease Progression in Friedreich Ataxia}
\begin{itemize}
\item Investigating the \textbf{blood transcriptomic landscape of Friedreich Ataxia (FRDA)} to identify molecular signatures associated with disease state and clinical severity.
\item Analyzing RNA-seq data from FRDA patients and controls to identify \textbf{differentially expressed genes and dysregulated biological pathways} associated with disease progression.
\item Integrating transcriptomic profiles with clinical severity measures to investigate relationships between molecular changes and the \textbf{phenotypic heterogeneity of FRDA}.
\end{itemize}

\project{iPSC-derived Cardiomyocyte Transcriptomics in Friedreich Ataxia}
\begin{itemize}
\item Investigating the molecular basis of \textbf{FRDA-associated cardiomyopathy} using iPSC-derived cardiomyocyte models, addressing a major disease manifestation and cause of mortality in FRDA.
\item Analyzing transcriptomic profiles across \textbf{Day 15 and Day 52 cardiomyocyte differentiation stages} to characterize transcriptional remodeling during cardiomyocyte maturation and disease progression.
\item Comparing cardiomyopathic and non-cardiomyopathic FRDA models to identify molecular pathways associated with the transition from \textbf{FXN deficiency to cardiac pathology}.
\item Investigating stage-specific transcriptional programs to identify candidate pathways associated with \textbf{cardiac dysfunction and disease heterogeneity}.
\end{itemize}

\project{Research Mentorship}
\begin{itemize}
\item Formally mentored a BTech student on the research project \textbf{``Haplotype Diversity and Shared Ancestry of the Sickle Cell Mutation (rs334) in Global Populations''}, providing guidance on population-genetic analysis, haplotype diversity, shared ancestry and interpretation of global allele-frequency patterns.
\item Guided the student through computational analysis, interpretation of population-genetic results and scientific presentation of the research findings.
\end{itemize}

%----------------------------------------------------------------------------------
% MSc THESIS
%----------------------------------------------------------------------------------

\entry
{MSc Thesis}
{2023--2024}
{National Centre for Cell Science (NCCS), Pune}
{}

\project{Computational Analysis of Viral Protein Evolution}
\begin{itemize}
\item Investigated \textbf{positive selection in viral proteins} to identify amino-acid substitutions associated with evolution and potential functional adaptation.
\item Developed custom amino-acid substitution matrices based on the \textbf{PAM framework} and mapped candidate substitutions onto protein structures to investigate their functional context.
\end{itemize}

%----------------------------------------------------------------------------------
% SUMMER RESEARCH
%----------------------------------------------------------------------------------

\entry
{Summer Research}
{2023}
{National Centre for Cell Science (NCCS), Pune}
{}

\project{Population-associated SARS-CoV-2 Genetic Variation}
\begin{itemize}
\item Investigated population-associated SARS-CoV-2 variants and their potential structural and functional consequences using sequence analysis and \textbf{molecular docking}.
\end{itemize}

%==================================================================================
% PUBLICATIONS
%==================================================================================

\section{Publications}

\begin{enumerate}
\item Subramanian, K., Bhattacharyya, C., \textbf{Mukherjee, A\textsuperscript{*}}, \textit{et al.} (2026). \textit{An Atlas of Indian Genetic Diversity}. \textbf{medRxiv}. Preprint. \href{https://doi.org/10.64898/2026.03.20.26348801}{DOI: 10.64898/2026.03.20.26348801}

\item Mondal, D., Bhattacharyya, C., \textbf{Mukherjee, A.}, \textit{et al.} (2026). \textit{Capturing India’s phenotypic diversity: Health insights from the GenomeIndia project}. \textbf{medRxiv}. Preprint. \href{https://doi.org/10.64898/2026.04.01.26349926}{DOI: 10.64898/2026.04.01.26349926}

\item Bhattacharyya, C., Subramanian, K., Uppili, B., \textbf{Mukherjee, A.}, \textit{et al.} (2025). \textit{Mapping Genetic Diversity with the GenomeIndia Project}. \textbf{Nature Genetics}, 57, 767--773. \href{https://doi.org/10.1038/s41588-025-02153-x}{DOI: 10.1038/s41588-025-02153-x}
\end{enumerate}

\textit{\textsuperscript{*}Equal first authorship. Ankit Mukherjee is bolded throughout.}
\vspace{0.4em}

%==================================================================================
% AWARDS
%==================================================================================

\section{Awards \& Academic Distinctions}

\begin{itemize}
\item \textbf{First Prize -- Best e-Poster Presentation}, Indian Ataxia Conclave 2026, for research on transcriptomic signatures and disease progression in Friedreich Ataxia.
\item \textbf{DBT-BET 2024 -- Category II}, qualified with 214 marks.
\item \textbf{All India Rank 51 -- GAT-B 2022}, Graduate Aptitude Test in Biotechnology.
\item \textbf{All India Rank 192 -- IIT-JAM Biotechnology 2022}.
\item \textbf{Qualified -- TIFR Graduate School Examination 2022}, Biology stream.
\item \textbf{Best Student Award -- BSc Microbiology}, University of Calcutta.
\end{itemize}

%==================================================================================
% TECHNICAL EXPERTISE
%==================================================================================

\section{Technical Expertise}

\skill{Programming \& Statistical Computing}{Python (NumPy, Pandas) \textbullet\ R (Tidyverse, Bioconductor, ggplot2) \textbullet\ Bash/Shell (awk, sed) \textbullet\ Perl}
\skill{Population \& Medical Genomics}{Population genomics \textbullet\ Pharmacogenomics \textbullet\ GWAS \textbullet\ Variant prioritization \textbullet\ Functional annotation}
\skill{Transcriptomics \& Functional Genomics}{Bulk RNA-seq analysis \textbullet\ Single-cell RNA-seq analysis \textbullet\ Gene-set and pathway enrichment \textbullet\ Gene co-expression analysis (WGCNA)}
\skill{Computational Infrastructure}{HPC \textbullet\ SLURM \textbullet\ Git/GitHub \textbullet\ Conda/Bioconda \textbullet\ Docker \textbullet\ Singularity}

%==================================================================================
% CONFERENCES & PRESENTATIONS
%==================================================================================

\section{Conferences \& Presentations}

\begin{itemize}
\item \textbf{Genomics India Conference 2026}. Presented a poster titled \textit{``Landscape of Clinically Actionable Pharmacogenomic Variation Across Diverse Indian Populations from the GenomeIndia Project.''}
\item \textbf{Workshop on AI-driven Multi-Omics}, jointly organized by \textbf{Indian Institute of Science (IISc)} and \textbf{IIT Madras}, held in conjunction with the Genomics India Conference 2026, Bengaluru.
\item \textbf{Indian Ataxia Conclave 2026}. Presented research on \textbf{blood transcriptomics and molecular signatures associated with Friedreich Ataxia disease severity}; awarded the \textbf{First Prize for Best e-Poster Presentation}.
\end{itemize}

%==================================================================================
% LANGUAGES
%==================================================================================

\section{Languages}

\textbf{English:} Professional proficiency \quad \textbf{Hindi:} Professional proficiency \quad \textbf{Bengali:} Native

\end{document}